\section{Obligation as Physical Weight}

A traditional Runekeeper tracks Obligation on a ledger. Numbers in a column. Segments on a clock. The debt is abstract, measurable, and—when paid—forgettable.

A Runic Warrior cannot forget.

Their Obligation is written on their body. A scar for a Rite invoked in desperation. A limp for a vow that cost more than expected. A missing finger for a debt that could not be deferred.

\subsection{The Scar as Receipt}

When a Runic Warrior marks Obligation, they choose where the cost appears.

\begin{itemize}
\item A \textbf{scar} on the arm or face (visible, memorable).
\item A \textbf{brand} on the back of the hand (visible to everyone they meet).
\item A \textbf{tattoo} that darkens or spreads (a record of accumulated debt).
\item A \textbf{crack} in their shield or gorget (if their Codex is stone or steel).
\item A \textbf{dullness} in their sentient blade's edge (the weapon shares the burden).
\end{itemize}

The player describes the mark. The GM may use it later as a narrative hook—an Inquisitor who recognizes the brand, a hedge-witch who reads the scars, a spirit who asks how the wound was earned.

\subsubsection{Mechanical Effects of Visible Obligation}

Obligation that is visibly marked on the body has both costs and benefits.

\textbf{Costs:}
\begin{itemize}
\item \textit{Social Penalty:} When dealing with authorities who oppose Runekeepers (Chain-Lanterns, Inquisitors, certain hedge-courts), the character suffers -1 die to social rolls. The marks announce their covenant.
\item \textit{Recognition:} Former enemies, rivals, or creditors may recognize the marks and seek the character out.
\item \textit{Healing Complications:} Magical healing that closes wounds does not erase Obligation scars. They remain, a permanent record.
\end{itemize}

\textbf{Benefits:}
\begin{itemize}
\item \textit{Intimidation:} When dealing with those who respect visible covenant marks (certain warrior cultures, mercenary companies, the Cord), the character gains +1 die to social rolls. The marks prove they have paid.
\item \textit{Witness:} The scars themselves count as a witness for Rites that require one. The body remembers what was sworn.
\item \textit{No Forgery:} A scar cannot be counterfeited. A brand cannot be stolen. The Obligation is undeniable.
\end{itemize}

\subsection{Fatigue as a Resource, Not a Penalty}

In the core rules, Fatigue is a penalty. It accumulates. It worsens Position. It leads to Harm.

For a Runic Warrior, Fatigue is also a resource.

A warrior who has been pushed to exhaustion knows exactly where their limit lies—and how to step just short of it. A body that has endured pain learns to read the shape of incoming harm.

\subsubsection{Talents That Use Fatigue}

This expansion introduces talents that convert Fatigue into advantage:

\begin{itemize}
\item \textit{Iron Resolve} (Advanced): When you would be reduced to 0 Harm, you may instead remain standing with 1 Harm and mark 2 Fatigue.
\item \textit{Scarred Resilience} (Foundation): Ignore 1 Fatigue from physical exertion once per scene.
\item \textit{Shared Burden} (Advanced): Transfer 1 Fatigue to your Thiasos. It becomes Winded (-1 die to all actions) for the rest of the scene.
\item \textit{Witnessed Scar} (Advanced): Carve a scar for each permanent Obligation. Once per session, touch a scar and gain +1 die on a roll directly tied to that Obligation.
\end{itemize}

\subsection{The Weight of Visibility}

A traditional Runekeeper can hide. Their Codex fits in a satchel. Their Obligation is ink on a page. They walk through a city gate without announcing their covenant.

A Runic Warrior cannot.

The scars are visible. The brand is undeniable. The cracked shield tells a story. The sentient blade hums when Inquisitors pass.

This visibility is a cost, but it is also a privilege. A Runic Warrior does not need to produce a writ. Their body is the writ. Their scars are the seal.

\subsection{Obligation and the Thiasos}

When a Runic Warrior marks Obligation, the Thiasos may also bear a mark.

\begin{itemize}
\item A warhorse develops a limp.
\item A spirit-hound's glow dims.
\item A sword's edge dulls.
\item A hollowed bird's milk-white eyes film over further.
\end{itemize}

The Runekeeper may choose to transfer the visible mark to their Thiasos (Sharing the Burden), but the Thiasos becomes Winded for the rest of the scene. Some Runekeepers refuse to burden their allies. Others consider it a fair distribution of weight.

\subsection{Paying Obligation}

A Runic Warrior pays Obligation in the same ways as any Runekeeper: service to their Patron, a witnessed atonement, or a sacrifice of memory or name.

But the payment is also \textit{visible}.

When a debt is cleared, the scar fades. The brand lightens. The crack in the shield seals. The sword's edge brightens.

This visibility is a comfort. The Runic Warrior can see when they are free. They do not need to consult a ledger. They only need to look at their hands.

\subsection{The Body as Ledger}

A Runic Warrior's body is a public record.

Every scar is a Rite invoked. Every limp is a battle survived. Every brand is a covenant declared.

This is not a metaphor. It is the central philosophy of the Runic Warrior: \textit{the debt is carried in the flesh, and the flesh does not forget.}

The thread held. The water carried. That is not a Rite. That is grace.
